\section{Relative \texorpdfstring{$\mathrm{TC}$}{TC} and integral lifts}\label{app:second-tc}
The proof of Theorem~\ref{thm:second-two-wild} compares a class in $K_2$ of a maximal order with $\mathrm{TC}$ at a finite truncation and then represents it by an integral class whose image in $K_2(D)$ is controlled; the proof of Theorem~\ref{thm:second-three-coker} uses the results on deep principal units and relative $K_1$.
Sections~\ref{app:second-tc-supports} and~\ref{app:second-tc-separation} provide the trace comparison at finite truncations and show that the image in $K_2(D)$ of a deep congruence kernel consists of transfers from a fixed splitting field; the remaining subsections show that, in a fixed truncation, images coming from sufficiently deep truncations are represented by integral classes.

\subsection{Supports and congruence towers}\label{app:second-tc-supports}
At finite truncations of a maximal order, we use the relative cyclotomic trace to compare integral relative $K$-classes.

\begin{lemma}\label{lem:app-nilpotent-trace}
Let $\Lambda$ be a unital associative ring in which $p$ is nilpotent, and let $I\subset\Lambda$ be a nilpotent two-sided ideal.
Then each $K_i(\Lambda,I)$ is a $p$-primary torsion group of bounded exponent, with a bound depending on $i$, and the homotopy groups of $K(\Lambda,I)$ are unchanged by $p$-completion.
The relative trace
\[
 K(\Lambda,I)^\wedge_p \xrightarrow{\ \sim\ } \mathrm{TC}(\Lambda,I)^\wedge_p
\]
is an equivalence; in particular $K_i(\Lambda,I)\simeq\pi_i\mathrm{TC}(\Lambda,I)^\wedge_p$ for all $i$.
\end{lemma}
\begin{proof}
The first assertion is Theorem A of Geisser--Hesselholt \cite{GeisserHesselholtNilpotent}, and the second follows from it.
The equivalence is the relative trace comparison for nilpotent ideals \cite{DundasGoodwillieMcCarthy}.
\end{proof}

Let $F$ be a complete discretely valued field, let $D$ be a central division algebra over $F$, and let $P=\mathcal O_D$ be its maximal order; let $J$ be the Jacobson radical of $P$, and let $\Delta=P/J$ be the residue division algebra of $D$.
Localization and d\'evissage give the fiber sequence
\[
 K(\Delta)\longrightarrow K(P)\longrightarrow K(D).
\]
Let $\pi$ be a uniformizer of the center, and put $P_s=P/\pi^sP$ and $J_s=J/\pi^sP$.
Then $J_s$ is nilpotent and $P_s/J_s=\Delta$.
If $p$ is the residue characteristic, then $p$ is nilpotent in $P_s$, so Lemma~\ref{lem:app-nilpotent-trace} applies to $(P_s,J_s)$, as well as to Iwahori and commutative truncations, and identifies relative classes at a fixed truncation integrally.
This comparison cannot be applied directly to the kernel $\pi^sP$ of $P\to P_s$.
To compare relative terms in the congruence tower with terms with supports, we use their respective continuity results \cite{DundasGoodwillieMcCarthy,GeisserHesselholtNilpotent,DundasMorrow,ClausenMathewMorrow}.

\begin{lemma}\label{lem:app-iwahori-derived-cut}
Let $R$ be a discrete valuation ring with uniformizer $\pi$, let $C$ be a maximal order finite free over $R$, and put $\Delta=C/J(C)$.
Let $H$ be the $r\times r$ Iwahori order with $C$ on and above the diagonal and $J(C)$ below it.
Let $J_H=J(H)$, and let $f$ be the last diagonal idempotent.
Let $H_s=H/\pi^sH$ and $C_s=C/\pi^sC$.
The transition map $Z_s\to Z_j$ between the relative corner cofibers
\[
 Z_s=\cofib\bigl(K(C_s,\Delta)
       \longrightarrow K(H_s,H/J_H)\bigr)
\]
is nullhomotopic whenever $s>j$.
Before truncation, the relative corner map $K(C,\Delta)\to K(H,H/J_H)$ is an equivalence.
\end{lemma}
\begin{proof}
The proof uses only that $C$ is finite free over $R$ and that $\pi C\subset J(C)$.
In $Z_s$, the map $K(C_s)\to K(H_s)$ is $-\otimes_{C_s}fH_s$ for right modules, and $K(\Delta)\to K(\Delta^r)$ is the inclusion of the last factor, where $H/J_H=\Delta^r$.
First consider the orders before truncation.
The right $C$-module $Hf$ over the corner $fHf=C$ is free, so $Hf\otimes^{\mathbb L}_CfH=Hf\otimes_CfH$, and multiplication $Hf\otimes_C fH\to H$ is injective.
The cokernel $\mathcal E$ of this multiplication map is the upper triangular algebra over $\Delta$ on the remaining $r-1$ vertices, so the derived quotient $H/^{\mathbb L}(HfH)$ by $f$ is equivalent to $\mathcal E$.
By its universal property of making $f$ a zero object, the derived quotient commutes with derived base change along the center.
Hence
\[
 H_s/^{\mathbb L}(H_sfH_s)
 \simeq\mathcal E_s,\qquad\mathcal E_s=\mathcal E\otimes_R^{\mathbb L}R/\pi^sR.
\]
Since $H,C$ are flat over $R$, the corner on the left is the actual ring $C_s$.
For a differential graded algebra $\mathcal E'$, such as $\mathcal E$ and $\mathcal E_s$, we write $K(\mathcal E')$ for the $K$-theory of the category $\operatorname{Perf}(\mathcal E')$ of perfect $\mathcal E'$-modules.
This identification is compatible with the transition maps \cite{BraunChuangLazarev}.

The two-term Koszul model of $R/\pi^sR$ gives $\mathcal E\oplus\mathcal E\epsilon_s$ as an $\mathcal E$-algebra model of $\mathcal E_s$, with $|\epsilon_s|=1$, $\epsilon_s^2=0$, and $d\epsilon_s=0$.
For $s\ge j$, the transition is given by $\epsilon_s\mapsto\pi^{s-j}\epsilon_j$.
Since $\pi\mathcal E=0$, for $s>j$ the transition factors as
\[
 \mathcal E_s\longrightarrow\mathcal E\longrightarrow\mathcal E_j,
\]
where the first map is the augmentation $\epsilon_s\mapsto0$ and the second is the inclusion in degree zero, a section of the augmentation of $\mathcal E_j$.
Thus $K(\mathcal E_s,\mathcal E)\to K(\mathcal E_j,\mathcal E)$ is nullhomotopic.

We compare the localization sequence for the derived quotient with the corner localization sequence for the residue ring $H/J_H=\Delta^r$.
The $K$-theory of an upper triangular algebra is the $K$-theory of its diagonal, so $K(\mathcal E)\simeq\bigoplus_{i=1}^{r-1}K(\Delta)$.
Placing $K(C_s)\to K(H_s)\to K(\mathcal E_s)$ over $K(\Delta)\to K(\Delta^r)\to K(\Delta^{r-1})$, with the right vertical map the augmentation followed by the diagonal projection, the vertical fibers give $Z_s\simeq K(\mathcal E_s,\mathcal E)$ compatibly with the transition maps.
This gives the required nullhomotopy.
Before truncation, the same comparison has zero cofiber.
\end{proof}

\begin{lemma}
\label{lem:app-nilpotent-etale-descent}
Let $R$ be a commutative Henselian local ring with residue field $k$ of characteristic $p$, let $\Lambda$ be a unital associative $R$-algebra, and let $I\subset\Lambda$ be a nilpotent ideal.
Let $a\ge1$.
For a finite \'etale $k$-algebra $U$, let $R_U$ be its Henselian lift, and write $\Lambda_U=\Lambda\otimes_RR_U$ and $I_U=I\otimes_RR_U$.
Let $\mathcal X$ be the sheaf of spectra defined by finite \'etale base change from
\[
 U\longmapsto K(\Lambda_U,I_U)/p^a.
\]
In particular, $\mathcal X(k)=K(\Lambda,I)/p^a$.
Let $\mathcal X_{\bar k}$ be the geometric stalk of $\mathcal X$.
Then there is a natural short exact sequence
\[
 0\longrightarrow H^1(k,\pi_3\mathcal X_{\bar k})
 \longrightarrow\pi_2\mathcal X(k)
 \longrightarrow(\pi_2\mathcal X_{\bar k})^{G_k}
 \longrightarrow0.
\]
\end{lemma}
\begin{proof}
Hesselholt--Larsen--Lindenstrauss prove \'etale base change for an arbitrary associative $R$-algebra \cite[Proposition 3.1, Addendum 3.2]{HesselholtLarsenLindenstrauss2020}.
Their Corollary~3.3 gives, for a finite Galois \'etale extension $R\to R'$ with group $G$,
\[
 \mathrm{TR}^n_i(\Lambda;p)\simeq
 H^0(G,\mathrm{TR}^n_i(\Lambda_{R'};p)),\qquad
 H^{q}(G,\mathrm{TR}^n_i(\Lambda_{R'};p))=0\quad(q>0).
\]
These results of Hesselholt--Larsen--Lindenstrauss assume neither that the residue field is finite nor that $p\ne2$.
We compare finite Postnikov truncations and then take the inverse limit to obtain $\mathrm{TR}^n(\Lambda;p)\simeq\mathrm{TR}^n(\Lambda_{R'};p)^{hG}$.
Taking the limits along restriction and Frobenius, we obtain $\mathrm{TC}(\Lambda;p)\simeq\mathrm{TC}(\Lambda_{R'};p)^{hG}$.
Since homotopy fixed points preserve finite limits and finite colimits of spectra, this descent also holds for relative terms and Moore coefficients.

Since $I$ is nilpotent, the relative trace comparison gives $K(\Lambda,I)/p^a\simeq\mathrm{TC}(\Lambda,I)^\wedge_p/p^a$, so the finite \'etale descent above also holds for $\mathcal X$.
Commutation of $\mathrm{TC}$ modulo $p^a$ with filtered colimits implies that $\mathcal X_{\bar k}\simeq\operatorname*{colim}_{k'}\mathcal X(k')$ and that each homotopy group of $\mathcal X_{\bar k}$ is a discrete $G_k$-module, where $k'$ runs over the finite separable extensions of $k$ in $\bar k$ \cite{DundasGoodwillieMcCarthy}, \cite[Theorem 2.7]{ClausenMathewMorrow}.
The homotopy groups of $\mathcal X$ and of $\mathcal X_{\bar k}$ are killed by $p^{2a}$, being extensions of $\pi_jK/p^a$ by $\pi_{j-1}K[p^a]$; in particular they are $p$-primary.
By \cite[Corollary 3.3]{HesselholtLarsenLindenstrauss2020} the higher group cohomology vanishes; the Eilenberg--Mac Lane sheaves of each $\mathrm{TR}^n$ are sheaves, so finite Postnikov truncations are hypersheaves, and connectivity, limits defining $\mathrm{TC}$, relative fibers and Moore coefficients preserve this.
Since $k$ has characteristic $p$, $\cd_p(k)\le1$ \cite[II, \S2.2, Proposition 3]{SerreGalois}, so the descent spectral sequence for each finite Postnikov truncation has only cohomological degrees zero and one.
The truncated tails lose at most one degree of connectivity under descent, so their connectivity still tends to infinity, and the continuous descent spectral sequence converges to give the displayed short exact sequence.
\end{proof}

\subsection{Congruence separation for a fixed splitting field}\label{app:second-tc-separation}
Next we show that the image in $K_2(D)$ of a sufficiently deep congruence kernel of $K_2(P)$ consists of transfers from a fixed splitting field.
This allows us to work at a finite truncation in the main text.

\begin{lemma}\label{lem:app-congruence-separation}
Let $F$ be a complete discretely valued field, let $D/F$ be a central division algebra, and let $M/F$ be a finite separable splitting field.
For all sufficiently large $s$,
\[
 \im\bigl(\ker[K_2(P)\to K_2(P_s)]\longrightarrow K_2(D)\bigr)
 \ \subseteq\ \theta_M K_2(M).
\]
\end{lemma}

\begin{proof}
The local rings $P$ and $P_s$ have Bass stable rank $1$.
Indeed, at least one entry of a left-unimodular pair is a unit; if the first entry is not a unit, adding the second makes it a unit.
The noncommutative stability theorem of Suslin--Tulenbaev therefore allows us to fix a single $n\ge5$, independent of $s$, with $K_2(P)\simeq\ker(\operatorname{St}_n(P)\to E_n(P))$ and with the analogous isomorphisms for $P_s$ and $D$ \cite{SuslinTulenbaev}.

Enlarging $n$ if necessary, which preserves these stability isomorphisms, we may assume that Corollary~\ref{cor:paper-congruence-splitting} applies with $E=M$.
Since stabilization identifies the kernels of both $\operatorname{St}_n(D)\to E_n(D)$ and the stable universal central extension with $K_2(D)$, the group $\operatorname{St}_n(D)$ is the pullback of the stable extension to $E_n(D)$.
Put
\[
 \Gamma_m=E_n(D)\cap(1+\pi^m\operatorname{Mat}_n(P)).
\]
By Corollary~\ref{cor:paper-congruence-splitting}, for sufficiently large $m$ there is a group-homomorphic section $\sigma$ over $\Gamma_m$ of $\operatorname{St}_n(D)/\theta_MK_2(M)\to E_n(D)$.

We take $s\ge 2m$.
By the presentation of the Steinberg group, the kernel of $\operatorname{St}_n(P)\to\operatorname{St}_n(P_s)$ is the normal closure of the root elements $x_{ij}(t)$ with $t\in\pi^sP$.
Indeed, in the quotient by this normal closure the coefficients can be written as elements of $P_s$, and the three types of Steinberg relations over $P_s$ can be checked using arbitrary lifts.
For distinct $i,j,k$,
\[
 x_{ij}(t)=[x_{ik}(\pi^m),x_{kj}(\pi^{-m}t)].
\]
Let $\gamma\in\operatorname{St}_n(P)$, and let $u_1$ and $u_2$ be the matrix images of $\gamma x_{ik}(\pi^m)\gamma^{-1}$ and $\gamma x_{kj}(\pi^{-m}t)\gamma^{-1}$.
Since $s\ge2m$, the matrix images of $x_{ik}(\pi^m)$ and $x_{kj}(\pi^{-m}t)$ lie in $\Gamma_m$, and conjugation by the image of $\gamma$ preserves $E_n(D)\cap(1+\pi^m\operatorname{Mat}_n(P))$, so $u_1,u_2\in\Gamma_m$.
In a central extension, changing either lift by a central element leaves the commutator unchanged.
Therefore the image of $\gamma x_{ij}(t)\gamma^{-1}$ in $\operatorname{St}_n(D)/\theta_MK_2(M)$ is $[\sigma(u_1),\sigma(u_2)]\in\sigma(\Gamma_m)$.

Every element of $\ker[K_2(P)\to K_2(P_s)]$ belongs to this normal closure and has trivial matrix image.
Its image in $\operatorname{St}_n(D)/\theta_MK_2(M)$ therefore lies in $\sigma(\Gamma_m)$ and maps to the identity in $E_n(D)$.
By the section property this image is trivial, so the generic-fiber image of the element lies in $\theta_MK_2(M)$.
\end{proof}

\subsection{Deep principal units and relative classes at finite truncations}\label{app:second-deep-units}
Lemma~\ref{lem:app-deep-unit-norm} controls the reduced norm on deep principal units.

For relative classes $y,z$, the relative trace comparison for nilpotent ideals (Lemma~\ref{lem:app-nilpotent-trace}) shows that equality of their images in relative $\mathrm{TC}^\wedge_p$ at a fixed truncation implies equality of their integral $K_2$ images at that truncation.
For corrections from commutative subfields and degree-zero cyclic orbits, we use Lemma~\ref{lem:app-orbit-stable-image} and Corollary~\ref{cor:app-commutative-cut-representative} below.

\begin{lemma}\label{lem:app-deep-unit-norm}
Let $T$ be a complete discretely valued field, let $B/T$ be a central division algebra, put $C=\mathcal O_B$, and let $\pi$ be a central element of positive valuation in $T$.
Let $\mathcal U_s(C)=\im(1+\pi^sC\to K_1(C))$.
There are integers $s_0$ and $t$ such that for all $s\ge s_0$ the reduced norm is injective on $\mathcal U_s(C)$ and
\[
 1+\pi^{s+t}\mathcal O_T\subseteq\Nrd_1(\mathcal U_s(C))
       \subseteq1+\pi^s\mathcal O_T.
\]
Moreover, for each $j$, every sufficiently deep unit of reduced norm one can be written as $\prod_i[b_i,c_i]$, with fixed $b_i\in C^\times$ and $c_i\in1+\pi^jC$.
These depths can be chosen uniformly over finite unramified extensions for which the base change of $C$ is the maximal order of a division algebra of the same degree.
When $B/T$ is a central simple algebra and $C\subset B$ is a finite $\mathcal O_T$-subalgebra satisfying $C\otimes_{\mathcal O_T}T=B$, we still have injectivity, the two inclusions for the norm image, and the commutator representation.
\end{lemma}
\begin{proof}
Fix integral units $b_i$ spanning $B$ over $T$.
The commutator product has differential $\sum_i(\operatorname{Ad}b_i-1)c_i$ at the identity.
By nondegeneracy of the reduced trace pairing, the orthogonal complement of the image of the differential is the common centralizer $T$ of the $b_i$, so the differential surjects onto $\ker\operatorname{Trd}_{B/T}$.
After clearing denominators, we apply the inverse function theorem with a fixed right inverse and realize every sufficiently deep element of reduced norm one using $c_i$ of any prescribed depth.
Since commutators vanish in $K_1(C)$, every sufficiently deep element of reduced norm one has trivial class in $K_1(C)$; this proves the injectivity of the reduced norm on $\mathcal U_s(C)$.
We obtain the left inclusion from Hensel's lemma in a direction with nonzero reduced trace, and the right inclusion from the reduced characteristic polynomial with integral coefficients.
Since unramified extension does not change the valuations of the denominators of the fixed Jacobian and of its right inverse, one depth works for all such base changes.
For a general finite subalgebra $C$, its units still span $B$ over $T$, since $1+\pi^ma\in C^\times$ for every lattice basis element $a$ and all sufficiently large $m$.
The same reduced trace pairing and fixed-Jacobian argument therefore applies.
\end{proof}

\subsection{Stable images of relative finite-group orbits}
At a fixed truncation, we show that the image on $\pi_2$ of the relative homotopy orbits of a sufficiently deep truncation agrees with that of the integral relative homotopy orbits.

Throughout this subsection, let $F$ be a complete discretely valued field, let $T/F$ be a finite Galois extension, put $G=\Gal(T/F)$, and let $\pi\in F$ have positive valuation.
Let $B/T$ be a central division algebra and put $C=\mathcal O_B$.
Let $H$ be either $C$ or the Iwahori order used in the main text with corner $C$, and write $J_H=J(H)$, $\bar H=H/J_H$, and $H_s=H/\pi^sH$.
We assume that $G$ acts semilinearly, preserving the order, its residue, and its truncations, and inducing the specified Galois action on the generic center $T$.
Through the corresponding action on the module category, we also use weak actions represented by ring automorphisms and coherence isomorphisms given by units.
For an Iwahori cover, we likewise use the actual action on the algebra formed from the direct sum of all vertices, on which vertex translation acts; in that case $H$ denotes this algebra.
In the homogeneous-unit type of the main text, identifying the vertices by powers of the unit $x$ gives $H\simeq M_r(C)$, and the weak action transported to $C$ by Morita equivalence falls under the case $H=C$.
This Morita equivalence is compatible with the residue and with each truncation, and it carries vertex translation to the weak action on the corner; at the generic point the cover side is $M_r(B)$ and the corner side is $B$.
For both types, $H_T=H\otimes_{\mathcal O_T}T\simeq M_r(B)$ at the generic point, and $N_{T/F}\Nrd_{H_T/T}$ is $G$-invariant on $K_1(H_T)$, so it sends the boundary used below, which vanishes in the coinvariants, to $1$.

\begin{lemma}
\label{lem:app-orbit-stable-image}
Let $Y=K(H,\bar H)$ and $Y_s=K(H_s,\bar H)$.
For every $j$ and all sufficiently large $s$,
\[
 \im[\pi_2((Y_s)_{hG})\to\pi_2((Y_j)_{hG})]
 =\im[\pi_2(Y_{hG})\to\pi_2((Y_j)_{hG})]
.
\]
When $G=1$, the same conclusion holds for the absolute spectra $Y=K(H)$ and $Y_s=K(H_s)$.

If the residue characteristic is $p$, then for every $j$ and every $x\in\pi_2((\mathrm{TC}(H,\bar H)^\wedge_p)_{hG})$ there is $y\in\pi_2(Y_{hG})$ such that $x$ and $\operatorname{tr}(y)$ have the same image in $\pi_2((\mathrm{TC}(H_j,\bar H)^\wedge_p)_{hG})\simeq\pi_2((Y_j)_{hG})$.
\end{lemma}

\begin{proof}
The order $H$ has global dimension one, and the rings $H_s$ and $\bar H$ are Artinian, so the negative $K$-groups of $H$, $H_s$, and $\bar H$ vanish \cite[III, \S4 and Exercise 4.3]{WeibelKbook}.
Since $H$ is complete, idempotents lift modulo $\pi^sH$ and modulo $J_H$, which gives the isomorphisms on $K_0$; together with surjectivity on $K_1$, they imply that $K(H,\pi^sH)$, $Y$, and $Y_s$ are concentrated in degrees at least one.
Thus, taking homotopy orbits of the $G$-equivariant fiber sequence $K(H,\pi^sH)\to Y\to Y_s$, we see that the boundary of each element of $\pi_2((Y_s)_{hG})$ belongs to $K_1(H,\pi^sH)_G$ and can be represented by $u\in1+\pi^sH$ \cite[III, Lemma 2.4(1)]{WeibelKbook}.
This boundary maps to zero in $(\pi_1Y)_G$, so the $G$-invariant norm $\lambda=N_{T/F}\circ\Nrd_{H_T/T}$ gives $\lambda(u)=1$.

We apply the argument of Lemma~\ref{lem:app-deep-unit-norm} over $F$.
We choose fixed $b_i\in H^\times$ spanning $H_T$ over $T$, and consider the map to $\ker\lambda$
\[
 \Phi((c_i),(a_\sigma))=\prod_i[b_i,c_i]\,
                   \prod_{\sigma\in G}\sigma(a_\sigma)a_\sigma^{-1}.
\]
Its differential is $\sum_i(\operatorname{Ad}b_i-1)x_i+\sum_\sigma(\sigma-1)y_\sigma$.
For the nondegenerate pairing $(x,y)\mapsto\Tr_{T/F}\operatorname{Trd}_{H_T/T}(xy)$, the orthogonal complement of the image of this differential is $T$ after the commutator terms, and then $T^G=F$ after the $G$-terms.
The image is therefore $\ker(\Tr_{T/F}\operatorname{Trd}_{H_T/T})$.
Since $d\lambda=\Tr_{T/F}\circ\operatorname{Trd}_{H_T/T}$ and both traces are surjective, $d\lambda$ is nonzero even when the characteristic divides the degree, and $\ker\lambda$ is smooth at the identity.
Using the same inverse function theorem with a fixed Jacobian, we see that, for every $j$ and all sufficiently large $s$, every $u\in1+\pi^sH$ with $\lambda(u)=1$ can be written as $u=\Phi((c_i),(a_\sigma))$ with $c_i,a_\sigma\in1+\pi^jH$.

Taking $b=b_i$ and $x=(c_i-1)b_i^{-1}$ in the relative elementary relation $[1+bx]=[1+xb]$ \cite[III, Exercise 2.4]{WeibelKbook} shows that $[b_i,c_i]$ vanishes in $K_1(H,\pi^jH)$.
The class of $\sigma(a_\sigma)a_\sigma^{-1}$ also vanishes in $K_1(H,\pi^jH)_G$, so the original boundary is zero at truncation $j$.
By exactness, the image at truncation $j$ of the given element of $\pi_2((Y_s)_{hG})$ lifts from the orbits of $Y$.
The depth depends only on $j$ and the fixed data.
The reverse inclusion holds because the map from $Y$ to $Y_j$ factors through $Y_s$.

For $G=1$, the same argument applies with $Y=K(H)$ and $Y_s=K(H_s)$, starting from $\Nrd_1(u)=1$, since the boundary maps to zero in $K_1(H)$.
The element $\pi$ may be any central element of positive valuation.

In residue characteristic $p$, the relative $K$-groups of a finite truncation are $p$-primary torsion groups of bounded exponent in each degree by Lemma~\ref{lem:app-nilpotent-trace}.
Only finitely many terms contribute in each orbit degree, so $p$-completion does not change the homotopy groups of $(Y_s)_{hG}$.
We apply the trace equivalence of the same lemma at a deep truncation, and then use the equality of images just proved to obtain an integral representative at the fixed truncation.
\end{proof}

\begin{corollary}
\label{cor:app-integral-cut-stable}
Let $T$ be a complete discretely valued field, let $B/T$ be a central division algebra, put $C=\mathcal O_B$ and $\Delta=C/J(C)$, and let $\pi\in T$ be any central element of positive valuation.
Let $C_s=C/\pi^sC$.
For every $j$ and all sufficiently large $s$,
\[
 \im[K_2(C_s)\to K_2(C_j)]
 =\im[K_2(C)\to K_2(C_j)].
\]
The same equality holds for the relative groups $K_2(C_s,\Delta)$ and $K_2(C,\Delta)$.
\end{corollary}
\begin{proof}
We apply the absolute and relative conclusions of Lemma~\ref{lem:app-orbit-stable-image} with $G=1$ and $H=C$.
\end{proof}

\subsection{Integral representatives over commutative valuation rings}
Finally, we record the commutative case of Lemma~\ref{lem:app-orbit-stable-image}, which is used for corrections by commutative subfields.

\begin{corollary}
\label{cor:app-commutative-cut-representative}
Let $R$ be a complete discrete valuation ring with residue field $\kappa$ of characteristic $p$.
For $j\ge1$, let $R_j$ be the quotient of $R$ by the $j$th power of a fixed uniformizer.
For every $j\ge1$ and every $x\in\pi_2\mathrm{TC}(R,\kappa)^\wedge_p$ there is $y\in K_2(R,\kappa)$ such that $x$ and $\operatorname{tr}(y)$ have the same image in $\pi_2\mathrm{TC}(R_j,\kappa)^\wedge_p\simeq K_2(R_j,\kappa)$.
\end{corollary}
\begin{proof}
We apply the last conclusion of Lemma~\ref{lem:app-orbit-stable-image} with $G=1$ and $H=R$.
\end{proof}
